\documentclass[12pt,a4paper]{article}
\usepackage[margin=25mm, headheight=15pt, headsep=10pt]{geometry}
\usepackage{fontspec}
\usepackage[table]{xcolor} % Note: table option is required for rowcolors
\usepackage{titlesec}
\usepackage{tocloft}
\usepackage{fancyhdr}
\usepackage{booktabs}
\usepackage{tcolorbox}
\tcbuselibrary{skins, breakable}
\usepackage[colorlinks=true, linkcolor=emerald, urlcolor=emerald, bookmarks=true]{hyperref}

\definecolor{emerald}{RGB}{0,102,51}
\definecolor{darkred}{RGB}{139,0,0}
\definecolor{lightgray}{RGB}{245,245,245}

\usepackage[nil]{babel}
\babelprovide[import, main]{bengali}
\babelprovide[import]{arabic}
\babelfont{rm}[Renderer=HarfBuzz,Scale=1.0]{Noto Serif Bengali}
\babelfont{sf}[Renderer=HarfBuzz]{Noto Sans Bengali}
\babelfont{tt}[Renderer=HarfBuzz]{Noto Sans Bengali}
\babelfont[arabic]{rm}[Renderer=HarfBuzz,Scale=1.1]{Noto Naskh Arabic}

% Section styling
\titleformat{\section}{\Large\bfseries\color{emerald}}{\thesection.}{0.5em}{}
\titleformat{\subsection}{\large\bfseries\color{emerald!85!black}}{\thesubsection.}{0.5em}{}
\titleformat{\subsubsection}{\normalsize\bfseries\color{emerald!70!black}}{\thesubsubsection.}{0.5em}{}
\titlespacing*{\section}{0pt}{18pt}{8pt}

% Fancy Header/Footer setup
\pagestyle{fancy}
\fancyhf{} % clear all header and footer fields
\fancyhead[L]{\color{emerald}\small\textbf{\nouppercase{\leftmark}}}
\fancyfoot[L]{\scriptsize\color{black!50}{\lat{AI}}-সংকলিত, আলেম-যাচাইকৃত নয়}
\fancyfoot[C]{\color{emerald!70!black}\thepage}
\renewcommand{\headrulewidth}{0.4pt}
\renewcommand{\headrule}{\hbox to\headwidth{\color{emerald}\leaders\hrule height \headrulewidth\hfill}}
\renewcommand{\footrulewidth}{0pt}

% Custom boxes
% 1. Ayat/Hadith Box
\newtcolorbox{ayatbox}[1][]{
    enhanced, breakable, 
    colback=emerald!5, 
    colframe=emerald, 
    boxrule=0.5pt, 
    arc=3pt, 
    left=10pt, right=10pt, top=10pt, bottom=10pt,
    #1
}

% 2. Quote/Translation Box
\newtcolorbox{quotebox}[1][]{
    enhanced, breakable,
    frame hidden,
    colback=lightgray,
    borderline west={3pt}{0pt}{emerald},
    left=15pt, right=10pt, top=10pt, bottom=10pt,
    sharp corners,
    #1
}

% 3. AI-generated notice box
\newtcolorbox{warnbox}[1][]{
    enhanced, breakable,
    colback=red!4,
    colframe=darkred,
    boxrule=0.8pt,
    arc=3pt,
    left=10pt, right=10pt, top=10pt, bottom=10pt,
    #1
}

% Commands
\newcommand{\arab}[1]{\foreignlanguage{arabic}{#1}}
\newcommand{\kib}[1]{\textcolor{emerald}{\textbf{#1}}}

% Latin (English) fallback: Noto Serif Bengali has NO Latin glyphs
\newfontfamily\latfont[Renderer=HarfBuzz]{Noto Serif}
\newcommand{\lat}[1]{{\latfont #1}}

\setlength{\parskip}{5pt}
\setlength{\parindent}{0pt}

% Fix for missing bullet in Noto Serif Bengali
\renewcommand{\labelitemi}{$\bullet$}



\begin{document}
\begin{center}{\Large\bfseries\color{emerald} আল্লাহর রহমত ও তাওবা — \lat{10}-হাতিব-ইবনে-আবি-বালতাআহ-ও-বদরিদের-ক্ষমা}\end{center}
\vspace{1em}
\section*{ঘটনা ১০ — হাতিব ইবনে আবি বালতাআহ: বড় ভুল, উমারের (রা.) ক্ষোভ ও নবীজির (সা.) ক্ষমার ঘোষণা}

\begin{warnbox}
 \textbf{গুরুত্বপূর্ণ নোটিশ (\lat{Important} \lat{Notice}):} এই কন্টেন্টটি \lat{AI} (কৃত্রিম বুদ্ধিমত্তা) দিয়ে প্রস্তুত ও সংকলিত। \lat{AI} কঠোর সূত্র-মান নীতি মেনে শুধু নির্ভরযোগ্য উৎস থেকে তথ্য সংগ্রহ করেছে — কেবল সহীহ/হাসান হাদিস ও সহীহ সনদের আসার রাখা হয়েছে, দুর্বল সনদ বাদ দেওয়া হয়েছে। তবে এটি কোনো আলেম বা মুহাদ্দিস কর্তৃক যাচাইকৃত নয়।
\end{warnbox}

\begin{quote}
\textbf{সম্পর্কিত ফাইল:} হাদিস ফাইল — ক্ষমার প্রশস্ততা; কুরআন ফাইল — ৪:৪৮-এর ব্যাখ্যা (শিরক ছাড়া সব পাপ আল্লাহর ইচ্ছাধীন)।
\end{quote}

\medskip\hrule\medskip

\subsection*{১. সূত্র কার্ড (\lat{Source} \lat{Card})}

\begin{table}[h]\centering\small
\begin{tabular}{ll}
বিষয় & বিবরণ \\
\hline
\textbf{ঘটনা} & মক্কা বিজয়ের প্রাক্কালে হাতিব (রাঃ) কুরাইশদের কাছে নবীজির (সা.) অভিযানের গোপন চিঠি পাঠালেন; চিঠি ধরা পড়লে উমার (রাঃ) মৃত্যুদণ্ড চাইলেন; নবীজি (সা.) বললেন — বদর যুদ্ধে অংশগ্রহণকারীদের আল্লাহ ক্ষমা করে দিয়েছেন \\
\textbf{বর্ণনাকারী} & আলী (রাঃ) — উবাইদুল্লাহ ইবনে রাফি'র সূত্রে \\
\textbf{গ্রন্থ ও নম্বর} & সহীহ বুখারি ৩০৮১; সহীহ মুসলিম ২৪৯৪ \\
\textbf{অধ্যায়} & বুখারি: কিতাবুল জিহাদ; মুসলিম: কিতাবু ফাদাইলিস সাহাবাহ — \lat{“}বদরিদের ফজিলত ও হাতিবের ঘটনা\lat{”} \\
\textbf{সনদের মান} & \textbf{সহীহ} (বুখারি-মুসলিম) \\
\textbf{স্থান ও সময়} & মদিনা — হিজরি ৮, মক্কা বিজয়ের আগে \\
\hline
\end{tabular}
\end{table}

\medskip\hrule\medskip

\subsection*{২. সংক্ষিপ্ত পটভূমি (\lat{Background})}

হাতিব ইবনে আবি বালতাআহ (রাঃ) ছিলেন \textbf{বদরের যুদ্ধে অংশগ্রহণকারী} সাহাবি — মুহাজিরদের মধ্যে একজন। তার পরিবার-পরিজন মক্কায় রয়ে গিয়েছিল, আর মুহাজিরদের অন্যান্যরা যাদের মক্কায় আত্মীয় ছিল, তারা পরিবারকে রক্ষার সুযোগ পেত। হাতিবের মক্কায় কেউ ছিল না — এই দুর্বলতা ও ভয় (\lat{fear} \lat{for} \lat{family}) তাকে একটি বড় ভুল করতে প্ররোচিত করল: তিনি নবীজির (সা.) মক্কা অভিযানের পরিকল্পনা কুরাইশদের জানিয়ে একটি গোপন চিঠি পাঠালেন। এটি ছিল গোপনীয়তার খিয়ানত (\lat{betrayal} \lat{of} \lat{trust}) — গুরুতর অপরাধ। ঘটনার পরিণতি ইসলামের ইতিহাসে রহমতের একটি অনন্য দৃষ্টান্ত।

\medskip\hrule\medskip

\subsection*{৩. পূর্ণ ঘটনার বিশুদ্ধ বর্ণনা (\lat{The} \lat{Full} \lat{Narration})}

আলী (রাঃ) বলেন:

\begin{quote}
রাসূলুল্লাহ (সা.) আমাকে, যুবাইরকে (রাঃ) এবং মিকদাদকে (রাঃ) পাঠিয়ে বললেন: \textbf{\lat{“}অমুক বাগানে যাও — সেখানে তোমরা এক মহিলাকে পাবে, যার কাছে হাতিব একটি চিঠি দিয়েছে; তা নিয়ে এসো।\lat{”}} আমরা সেখানে গিয়ে মহিলার কাছে চিঠি চাইলাম। সে বলল: \lat{“}আমার কাছে কোনো চিঠি নেই।\lat{”} আমরা বললাম: \lat{“}চিঠি বের করো — নইলে তোমার কাপড় খুলে ফেলব।\lat{”} তখন সে তার \textbf{চুলের বেণী থেকে চিঠিটি বের করল।}
\end{quote}

\begin{quote}
আমরা চিঠিটি রাসূলুল্লাহর (সা.) কাছে পৌঁছে দিলাম — তাতে হাতিব মক্কার মুশরিকদের কাছে নবীজির (সা.) অভিযানের কিছু খবর জানিয়েছিলেন। নবীজি (সা.) হাতিবকে ডেকে জিজ্ঞেস করলেন: \textbf{\lat{“}হাতিব! এটা কী?\lat{”}} তিনি বললেন: \lat{“}হে আল্লাহর রাসূল! আমার ব্যাপারে তাড়াহুড়া করবেন না। আল্লাহর কসম! আমি কুফরি করিনি এবং ইসলামের প্রতি ভালোবাসা ছাড়া আর কিছুই বাড়েনি। কিন্তু আপনার সাথে থাকা মুহাজিরদের কারো না কারো মক্কায় আত্মীয় আছে, যারা তাদের পরিবার-সম্পদ রক্ষা করে — আর আমার সেখানে কেউ নেই। তাই আমি তাদের কাছে একটি উপকারের আশা করলাম, যেন তারা আমার পরিবারকে রক্ষা করে। আমি কুফরি বা ধর্মত্যাগের জন্য এটা করিনি।\lat{”} নবীজি (সা.) তাকে সত্যবাদী জেনে গ্রহণ করলেন।
\end{quote}

\begin{quote}
তখন উমার (রাঃ) দাঁড়িয়ে বললেন: \textbf{\lat{“}হে আল্লাহর রাসূল! আমাকে অনুমতি দিন — আমি এই মুনাফিকের গর্দান উড়িয়ে দিই!\lat{”}} নবীজি (সা.) বললেন: \textbf{\lat{“}তুমি কী জানো? সম্ভবত আল্লাহ বদরের যুদ্ধে অংশগ্রহণকারীদের প্রতি দৃষ্টি দিয়েছেন এবং বলেছেন: 'তোমরা যা ইচ্ছা করো — আমি তোমাদের ক্ষমা করে দিয়েছি।'\lat{”}}
\end{quote}

(সহীহ বুখারি ৩০৮১; সহীহ মুসলিম ২৪৯৪)

এ ঘটনার পরই সূরা আল-মুমতাহিনাহর (৬০:১) প্রথম আয়াত নাজিল হয় — \lat{“}হে ঈমানদারগণ! তোমরা আমার শত্রু ও তোমাদের শত্রুকে বন্ধু বানাবে না…\lat{”} (মুসলিম ২৪৯৪-এর বর্ণনায়)।

\medskip\hrule\medskip

\subsection*{৪. শব্দের অর্থ (\lat{Key} \lat{Words})}

\begin{table}[h]\centering\small
\begin{tabular}{ll}
শব্দ & অর্থ \\
\hline
\textbf{খিয়ানাহ (\arab{خيانة})} & বিশ্বাস ভঙ্গ — আমানতের খিয়ানত \\
\textbf{নাফাকা (\arab{نافق})} & মুনাফিক \lat{behavior} প্রদর্শন করা — উমারের (রা.) অভিযোগের ভাষা \\
\textbf{আহলু বদর (\arab{أهل} \arab{بدر})} & বদর যুদ্ধে অংশগ্রহণকারী সাহাবিগণ \\
\textbf{ইত্তালাআ (\arab{اطّلع})} & তীক্ষ্ণ দৃষ্টি দেওয়া — আল্লাহর বিশেষ অনুগ্রহের দৃষ্টি \\
\hline
\end{tabular}
\end{table}

\medskip\hrule\medskip

\subsection*{৫. সংশ্লিষ্ট দলিল ও রেফারেন্স}

\begin{itemize}
\item \textbf{আয়াত (আল-আনফাল ৮:২৭):} \lat{“}তোমরা আল্লাহ ও রাসূলের খিয়ানত করো না…\lat{”} — হাতিবের (রাঃ) ভুলের প্রেক্ষাপটে অনেক তাফসীরকার এই আয়াতের আলোচনা করেছেন; তবে নাজিলের নির্দিষ্ট কারণ নিয়ে মতভেদ আছে (ইবনে কাসীর এখানে হাতিব-চিঠির ঘটনার বদলে বদর-ক্ষমার হাদিস উল্লেখ করেছেন)।
\item \textbf{আয়াত (আল-মুমতাহিনাহ ৬০:১):} শত্রুর সাথে গোপন সম্পর্ক নিষেধ — হাতিবের (রাঃ) ঘটনার পরই নাজিল (মুসলিম ২৪৯৪)।
\item \textbf{হাদিস (বুখারি ৩০৮১; মুসলিম ২৪৯৪):} \lat{“}সম্ভবত আল্লাহ বদরিদের দিকে দৃষ্টি দিয়ে বলেছেন: তোমরা যা চাও — আমি ক্ষমা করেছি।\lat{”}
\item \textbf{হাদিস (বুখারি ৩৯৮৩; মুসলিম ২৪৯৪-এর নিকট):} বদরিদের মর্যাদা — নবীজি (সা.) বলেন: \lat{“}আল্লাহ হয়তো বদরিদের দিকে তাকিয়ে বলেছেন…\lat{”}
\end{itemize}
\subsubsection*{মূল শিক্ষা (দলিলভিত্তিক)}

\begin{enumerate}
\item \textbf{বড় ভুলও রহমতের গণ্ডির বাইরে নয়:} হাতিব (রাঃ) — বদরের সাহাবি — একটি গুরুতর ভুল করলেন; উমার (রাঃ) মৃত্যুদণ্ড চাইলেন; কিন্তু নবীজি (সা.) জানালেন — তার অতীতের মহান আমল (বদর) আল্লাহর রহমতের দৃষ্টিতে আছে। \textbf{একটি ভুলে একজন মুমিনের পুরো অতীত মুছে যায় না।}
\item \textbf{উমার (রা.)-এর জবাব ও নবীজির (সা.) শিক্ষা:} উমার (রাঃ) খাঁটি ঈমান থেকে কঠোরতা চেয়েছিলেন; নবীজি (সা.) শিখালেন — সন্দেহ ও ক্ষোভের আগে আল্লাহর রহমতের সম্ভাবনার কথা ভাবতে হবে। কঠোরতা ও রহমতের ভারসাম্য (\lat{balance}) এখানেই।
\item \textbf{দুর্বলতা স্বীকার করা:} হাতিব (রাঃ) ভুল স্বীকার করে তার কারণ (পারিবারিক নিরাপত্তার ভয়) স্পষ্টভাবে বললেন — মিথ্যা বলেননি, অজুহাত দেখাননি। নবীজি (সা.) তাকে সত্যবাদী জেনে গ্রহণ করলেন। \textbf{পাপের পর সত্য বলা ও কারণ খোলা বলা — ক্ষমার পথ সহজ করে।}
\item \textbf{এই ঘটনা থেকে একজন পাপীর শিক্ষা:} অতীতের কোনো মহান আমল বা সৎ জীবন এখনো আল্লাহর হিসাবে মূল্যবান — ভুলের পর হতাশ হয়ে সবকিছু ছেড়ে দেওয়ার প্রয়োজন নেই; ফিরে এসে সত্য বলাই পথ।
\end{enumerate}
\medskip\hrule\medskip

\subsection*{৬. রেফারেন্স}

\begin{itemize}
\item সহীহ বুখারি ৩০৮১ (কিতাবুল জিহাদ ওয়াস সিয়ার)
\item সহীহ মুসলিম ২৪৯৪ (কিতাবু ফাদাইলিস সাহাবাহ)
\item কুরআন: আল-মুমতাহিনাহ ৬০:১; আল-আনফাল ৮:২৭
\item তাফসীর: ইবনে কাসীর (৮:২৭-এর আলোচনা)
\item ব্যাখ্যা: ফাতহুল বারি — হাতিবের ঘটনার ব্যাখ্যা
\end{itemize}

\end{document}
